\section{Overall Algorithm}
\label{sec:overall}

With the trained channel denoisers and activity classifiers, the
updates in Section~\ref{sec:algorithm} yield the unfolded receiver.
We describe its message scheduling, variance calibration and damping,
and computational complexity.

\subsection{Message Scheduling}

Given normalized $\mathbf Y$, the operator $\mathbf A_{\rm c}$,
noise variance $\sigma_{\rm c}^2$, and activity prior $\rho$,
initialize
\begin{equation}
 \mathbf M_1=\mathbf0,\qquad w_1=w_{\rm ref}=\rho P_H,
 \label{eq:twc-initialization}
\end{equation}
where $P_H=\E_{q_{\rm tr}}\|h\|^2/D$ is the average channel power
per coefficient. The linear estimator first forms the user queries,
which are processed by the channel denoisers and activity classifiers.
Their outputs determine the effective channel estimates and the
feedback to the next layer. After $L$ layers, the receiver returns
$\widehat{\mathbf X}_L$ and declares user $k$ active if $p_{kL}\geq1/2$.
The effective channel estimates are rescaled by $s$.
The complete procedure is summarized in Table~\ref{tab:twc-algorithm}.

\begin{table}[t]
\centering
\caption{Unfolded receiver for joint activity detection and channel estimation}
\label{tab:twc-algorithm}
\begin{tabular}{p{0.93\linewidth}}
\toprule
\textbf{Input:} $\mathbf Y,\mathbf A_{\rm c},\sigma_{\rm c}^2,\rho$,
trained layer parameters, and $P_H$.\\
\textbf{Initialize:} $(\mathbf M_1,w_1)$ by \eqref{eq:twc-initialization}.\\
\textbf{For} $\ell=1,\ldots,L$:\\
1. Compute the linear update and extrinsic query using
\eqref{eq:twc-linear-correction} and \eqref{eq:twc-query-calibrated}.\\
2. Estimate the conditional channels and activity probabilities
using \eqref{eq:twc-conditional-mean} and \eqref{eq:twc-activity}.\\
3. Form the effective channels using \eqref{eq:twc-estimate}.\\
4. If $\ell<L$, calculate the residual scale by \eqref{eq:twc-residual-scale},
calibrate the feedback variance by
\eqref{eq:twc-positive-scale}--\eqref{eq:twc-feedback-ratio}, and
form the extrinsic message by \eqref{eq:twc-feedback-proposal}.\\
5. If $\ell<L$, damp the outgoing message using
\eqref{eq:twc-damping-weight}--\eqref{eq:twc-damping-mean}.\\
\textbf{Return:} $\widehat{\mathbf X}_L$, $p_{kL}$,
and $\widehat\alpha_k=\mathbf1\{p_{kL}\geq1/2\}$.\\
\bottomrule
\end{tabular}
\end{table}

\subsection{Variance Calibration and Damping}
\label{sec:twc-damping}

The scalar residual in \eqref{eq:twc-residual-scale} is used to
construct the feedback variance. To obtain a positive scale, we
apply a smooth positive mapping and a trainable multiplier.
Let $\operatorname{softplus}(x)=\log(1+e^x)$ and define
\begin{align}
 T_\ell&=\eta v_\ell,\qquad
 z_\ell=e^{c_{\beta,\ell}}
 [T_\ell\operatorname{softplus}(e_\ell/T_\ell)+\epsilon_s],
 \label{eq:twc-positive-scale}\\
 q_\ell&=\epsilon_r+(q_{\max}-\epsilon_r)
 \tanh\!\left(\frac{z_\ell}{(q_{\max}-\epsilon_r)v_\ell}\right),
 \qquad \tau_\ell=v_\ell q_\ell.
 \label{eq:twc-feedback-ratio}
\end{align}
Here $\eta$ controls the mapping temperature, $\epsilon_s>0$
sets a positive floor, and $0<\epsilon_r<q_{\max}<1$ bounds
the variance ratio. Consequently, the feedback variance satisfies
$0<\tau_\ell<v_\ell$, as required by
\eqref{eq:twc-feedback-proposal}.

To smooth the updates between layers, we damp the Gaussian messages
in their natural parameters \cite[Sec.~5.2]{vehtari2020}:
\begin{align}
 \delta_\ell&=\min\{1,w_{\rm ref}/\widetilde w_{\ell+1}\},
 \label{eq:twc-damping-weight}\\
 w_{\ell+1}^{-1}&=(1-\delta_\ell)w_\ell^{-1}
                   +\delta_\ell\widetilde w_{\ell+1}^{-1},
 \label{eq:twc-damping-precision}\\
 \mathbf M_{\ell+1}&=w_{\ell+1}
 \left[(1-\delta_\ell)\frac{\mathbf M_\ell}{w_\ell}
       +\delta_\ell\frac{\widetilde{\mathbf M}_{\ell+1}}
                          {\widetilde w_{\ell+1}}\right].
 \label{eq:twc-damping-mean}
\end{align}
The weight $\delta_\ell$ decreases when the proposed variance exceeds
the initial reference $w_{\rm ref}$. Otherwise the proposed message is
passed directly to the next layer.

\subsection{Computational Complexity Analysis}

The main computations are the linear estimation and the forward
passes through the two networks. The linear update can be evaluated
using a single singular value decomposition,
$\mathbf A_{\rm c}=\mathbf U\mathbf S\mathbf V^{\mathsf H}$.
With $s_i$, $i=1,\ldots,r_A$, denoting its nonzero singular values,
\eqref{eq:twc-linear-correction} becomes
\begin{equation}
 \bm\Delta_\ell=\mathbf V
 \diag\left\{\frac{w_\ell s_i}{\sigma_{\rm c}^2+w_\ell s_i^2}\right\}
 \mathbf U^{\mathsf H}(\mathbf Y-\mathbf A_{\rm c}\mathbf M_\ell).
 \label{eq:twc-svd-correction}
\end{equation}
The variance fractions are evaluated as
\begin{equation}
 f_\ell=\frac{J-r_A+\sum_{i=1}^{r_A}
 \sigma_{\rm c}^2/(\sigma_{\rm c}^2+w_\ell s_i^2)}{J},\qquad
 d_\ell=\frac1J\sum_{i=1}^{r_A}
 \frac{w_\ell s_i^2}{\sigma_{\rm c}^2+w_\ell s_i^2}.
 \label{eq:twc-linear-fractions}
\end{equation}
The singular vectors and values are reused across all layers.

Let $r=\min(Q,J)$. A dense thin SVD has complexity $O(QJr)$.
In each layer, forming the observation residual costs $O(QJC)$,
and applying the singular vectors in \eqref{eq:twc-svd-correction}
costs $O((Q+J)rC)$. The feedback residual has the same order as the
observation residual. Denote the costs of a score-network and
activity-head evaluation for one user by $F_s$ and $F_a$,
respectively. The total complexity is therefore
\begin{equation}
 O\left(QJr+
 L\left[QJC+(Q+J)rC+K(F_s+F_a)\right]\right).
 \label{eq:twc-complexity}
\end{equation}
For $Q\leq J$, this reduces to
$O(Q^2J+L[QJC+K(F_s+F_a)])$.
If the score network and activity head contain $P_s$ and $P_a$
parameters, the $L$ layers store $L(P_s+P_a+3)$ parameter scalars.
The phase and noise estimators in Appendix~\ref{app:frontend}
are executed once before these layers.
